% !TEX root = Tesi_Corsi_Leonardo_656276.tex
\chapter{Implementazione}
\label{ch:implementazione}

\section{Le quattro istruzioni di comunicazione}
\label{sec:istruzioni}

L'estensione aggiunge quattro istruzioni, tutte con l'opcode \emph{custom-0}
descritto nella sezione~\ref{sec:custom0} e distinte dal campo \texttt{funct3}:

\begin{center}
\begin{tabular}{llll}
    \toprule
    istruzione & \texttt{funct3} & punta & restituisce \\
    \midrule
    \texttt{IN rd, dir}     & 0 & il canale del vicino & la parola letta \\
    \texttt{OUT rs, dir}    & 1 & il canale proprio & niente \\
    \texttt{ISRDY rd, dir}  & 2 & il canale del vicino & 1 se c'è un dato da leggere \\
    \texttt{SETRDY rd, dir} & 3 & il canale proprio & 1 se ha pubblicato \\
    \bottomrule
\end{tabular}
\end{center}

Le due istruzioni di test sono \textbf{non bloccanti}: non sospendono la cella,
ma restituiscono un esito in un registro, e spetta al programma decidere cosa
farne, di norma ripetendo il tentativo. La scelta è motivata dall'hardware. Un
test bloccante è una primitiva di sincronizzazione, e in un circuito richiede
uno stallo della pipeline oppure uno scheduler; un test non bloccante è un filo
letto come un bit, il cui costo implementativo è nullo e il cui costo in cicli
è pagato esplicitamente dal programma, dove resta visibile e misurabile.

\texttt{IN} consuma leggendo: preleva il dato e, se il canale era pieno, libera
lo slot. La lettura di un canale vuoto restituisce il valore precedente e non
altera lo stato; non costituisce un errore, ma presuppone che il programma
abbia verificato la disponibilità del dato con \texttt{ISRDY}.

\section{Il ciclo di clock a due fasi}
\label{sec:due-fasi}

Ogni campo del canale ha il suo gemello \texttt{\_next}, e vale una regola
unica: \textbf{le letture guardano sempre lo stato attuale, le scritture vanno
sempre nel gemello}. Il ciclo di clock è quindi diviso in due fasi:

\begin{lstlisting}[caption={Il ciclo di clock globale},label={lst:gridstep},language=C,float=htbp]
void grid_step(Grid *grid) {
    /* FASE COMPUTE: legge lo stato attuale,
       scrive solo nei campi _next */
    for (ogni cella)  execute_step(...);

    /* FASE COMMIT: i campi _next diventano attuali */
    for (ogni cella)
        for (i suoi 4 canali di uscita) ch_commit(...);
}
\end{lstlisting}

Senza la separazione, il comportamento dipenderebbe dall'ordine di visita
dell'array: una cella già eseguita vedrebbe il valore \emph{nuovo} del vicino,
una non ancora eseguita quello \emph{vecchio}, e lo stesso programma darebbe
risultati diversi al variare dell'ordine dei due cicli annidati. Con
l'aggiornamento separato la semantica è quella di un circuito sincrono reale:
tutte le celle leggono lo stato al fronte $k$ e scrivono lo stato al fronte
$k+1$.

La conseguenza è misurabile ed è ciò che rende il modello sistolico~\cite{KungLeiserson}:
\textbf{ogni salto verso un vicino costa esattamente un ciclo}. Nell'aggiornamento
di fine ciclo, riportato nel codice~\ref{lst:commit}, i due contatori passano
sempre, mentre il dato passa solo sulla transizione di \texttt{wp}.

\begin{lstlisting}[caption={L'aggiornamento di fine ciclo},label={lst:commit},language=C,float=htbp]
static inline void ch_commit(Channel *c) {
    /* il confronto E' il segnale di pubblicazione */
    if (c->wp_next != c->wp) {
        c->data = c->data_next;
    }
    c->wp = c->wp_next;
    c->rp = c->rp_next;
}
\end{lstlisting}

I due registri del dato non sono due buffer generici, ma \textbf{due registri
in serie}: il registro di uscita del mittente è \texttt{data\_next}, quello
letto dal vicino è \texttt{data}, e l'abilitazione del secondo è il fronte di
\texttt{wp}. In assenza di quella condizione, una \texttt{OUT} a canale pieno
cancellerebbe un valore già pubblicato e non ancora letto.

\section{Il protocollo}
\label{sec:protocollo}

La coppia \texttt{OUT} e \texttt{SETRDY} è asimmetrica, e l'asimmetria è
deliberata:

\begin{lstlisting}[caption={Le due metà della pubblicazione},label={lst:outsetrdy},language=C,float=htbp]
/* OUT: carica il registro di uscita. Non fallisce mai. */
static inline void ch_write(Channel *c, uint32_t v) {
    /* congelato se la pubblicazione di questo ciclo
       e' gia' decisa */
    if (c -> wp_next == c -> wp) {
        c -> data_next = v;
    }
}

/* SETRDY: l'unica che puo' essere rifiutata,
   e a rifiutarla e' il comparatore. */
static inline int ch_setrdy(Channel *c) {
    if (ch_iswrt(c) && c -> wp_next == c -> wp) {
        c -> wp_next = c -> wp ^ 1;
        return 1;
    }
    return 0;
}
\end{lstlisting}

\textbf{\texttt{OUT} non è un invio}, ma il caricamento di un registro di
uscita. Il canale fornisce una garanzia solo a partire da \texttt{SETRDY}, che è
l'istruzione con cui il produttore commuta il proprio \emph{ready bit}, e la
garanzia è che \emph{ogni valore pubblicato viene consegnato esattamente una
volta}. Due \texttt{OUT} consecutive senza una \texttt{SETRDY} interposta sono
due caricamenti dello stesso registro, non due messaggi: il primo valore viene
sovrascritto senza conseguenze, poiché non era ancora stato promesso ad alcun
consumatore. È la stessa distinzione che in un \emph{handshake}
hardware separa il registro di uscita dal segnale di \emph{valid}.

\subsection{La finestra fra \texttt{OUT} e \texttt{SETRDY}}

\texttt{OUT} e \texttt{SETRDY} sono due istruzioni in due cicli distinti, poiché
una cella esegue una istruzione per ciclo. Nell'intervallo il consumatore può
leggere, e da questa possibilità nasce l'unico caso delicato del protocollo:

\begin{center}
\begin{tabular}{ll}
    \toprule
    $t$ & \texttt{OUT s1, EST} $\rightarrow$ il canale è pieno \\
        & il consumatore legge, quindi al ciclo successivo è vuoto \\
    $t+1$ & \texttt{SETRDY t0, EST} $\rightarrow$ ora riesce, ma pubblica \emph{quale} valore? \\
    \bottomrule
\end{tabular}
\end{center}

Se il registro di uscita non fosse stato caricato al ciclo $t$, la
\texttt{SETRDY} consegnerebbe il valore vecchio, già consegnato una volta:
un duplicato, e il valore nuovo non partirebbe mai. Si tratterebbe di un errore
silenzioso.

La soluzione risiede in una \textbf{asimmetria deliberata}: la \texttt{OUT} non
è mai rifiutata. Carica sempre \texttt{data\_next}, anche a canale pieno, ed è
un'operazione innocua perché il consumatore osserva \texttt{data}. Al ciclo $t$
il valore nuovo è quindi già nel registro, e la \texttt{SETRDY} di $t+1$
consegna esattamente quel valore, senza che il canale debba registrare il
passaggio di una \texttt{OUT}.

\paragraph*{Invariante.} \texttt{data} cambia \textbf{solo} quando \texttt{wp}
commuta.

\paragraph*{Dimostrazione.} \texttt{data} è scritto in un solo punto, sotto la
condizione $\mathtt{wp\_next} \neq \mathtt{wp}$ dell'aggiornamento di fine
ciclo. Un valore pubblicato resta quindi intatto finché non viene pubblicato il
successivo, e la pubblicazione richiede $\mathtt{wp} = \mathtt{rp}$, ovvero che
il consumatore abbia letto. Nessuna \texttt{OUT}, per quante ne esegua il
produttore, può raggiungere \texttt{data} senza passare da una \texttt{SETRDY}
riuscita. $\square$

Ne segue che una \texttt{OUT} non può mai fallire e che un valore caricato non
va mai perso: al più viene sovrascritto da una \texttt{OUT} successiva prima di
essere pubblicato, che è precisamente la semantica voluta.

\subsection{Il ciclo completo}

Si consideri un produttore in $(0,0)$ che invia 42 e poi 43 a est, con il
consumatore in $(0,1)$:

\begin{center}
\begin{tabular}{clll}
    \toprule
    ciclo & produttore & consumatore & stato a fine ciclo \\
    \midrule
    0 & \texttt{OUT s1, EST(42)} & \texttt{ISRDY t0, OVEST} $\rightarrow$ 0 & vuoto, \texttt{data} intatto \\
    1 & \texttt{SETRDY t0, EST} $\rightarrow$ 1 & \texttt{ISRDY t0, OVEST} $\rightarrow$ 0 & \textbf{pieno}, \texttt{data} cattura 42 \\
    2 & (istruzione successiva) & \texttt{ISRDY t0, OVEST} $\rightarrow$ 1 & pieno \\
    3 & \texttt{OUT s1, EST} (43) & \texttt{IN t1, OVEST} $\rightarrow$ 42 & \textbf{vuoto}, \texttt{data} resta 42 \\
    4 & \texttt{SETRDY t0, EST} $\rightarrow$ 1 & \dots & pieno, \texttt{data} cattura 43 \\
    \bottomrule
\end{tabular}
\end{center}

Le righe rilevanti sono la 3 e la 4. Al ciclo 3 il produttore carica il valore
successivo \emph{mentre} il consumatore sta leggendo il precedente: la
\texttt{OUT} va a segno ugualmente ed è innocua, poiché tocca solo il registro
di uscita. Al ciclo 4 il canale si è svuotato e la \texttt{SETRDY} consegna 43,
il valore corretto, una sola volta e senza un giro di ritentativo.

\section{Stato dei registri}
\label{sec:registri}

Ogni cella conosce la propria posizione fin dall'avvio. Alla costruzione della
griglia, quattro registri argomento sono precaricati:

\begin{center}
\begin{tabular}{lll}
    \toprule
    registro & contenuto & \\
    \midrule
    \texttt{a0} & riga della cella & $r$ \\
    \texttt{a1} & colonna della cella & $c$ \\
    \texttt{a2} & righe totali & $R$ \\
    \texttt{a3} & colonne totali & $C$ \\
    \bottomrule
\end{tabular}
\end{center}

L'identità è precaricata anziché ricavata da un identificatore lineare per una
ragione precisa: il passaggio da $\mathit{id}$ a $(r,c)$ richiede una
\textbf{divisione}, e RV32I non dispone di istruzioni di divisione. Ogni cella
dovrebbe quindi eseguire all'avvio una routine software di decine di istruzioni
al solo scopo di determinare la propria posizione. Nel circuito reale la
posizione è cablata, dunque precaricarla è anche il modello più fedele.

È questo che abilita l'idea centrale dei programmi: \textbf{una sola immagine,
caricata identica in tutte le celle, che sceglie il proprio ruolo dalla
posizione}.

Il resto della convenzione sui registri è locale ai programmi e segue l'uso
tipico dell'architettura: i registri temporanei \texttt{t0}--\texttt{t2}
raccolgono gli esiti dei test, i registri salvati \texttt{s1}--\texttt{s2}
tengono l'accumulatore e il contatore del ciclo principale. Nei due programmi
usati per la taratura delle misure, \texttt{s3} e \texttt{s4} contano in
assembly le pubblicazioni rifiutate e le attese a vuoto, e servono a verificare
i contatori del simulatore (sezione~\ref{sec:taratura}). Una cella termina
eseguendo \texttt{ecall}, che la ferma senza fermare le altre.

\section{L'ingresso e l'uscita sul perimetro}
\label{sec:io-bordo}

Il programma ospite si affaccia sui canali di bordo con due sole funzioni, che
si limitano a ripetere ciò che farebbe una cella vicina:

\begin{lstlisting}[caption={L'ospite come vicino},label={lst:pushpop},language=C,float=htbp]
int grid_push(Grid *grid, int r, int c, int dir, uint32_t v) {
    Channel *ch = grid_at(grid, r, c) -> in_ch[dir];
    ch_write(ch, v);        /* la OUT che farebbe il vicino */
    return ch_setrdy(ch);   /* la sua SETRDY, stesso esito */
}

int grid_pop(Grid *grid, int r, int c, int dir, uint32_t *v) {
    Channel *ch = &grid_at(grid, r, c) -> out_ch[dir];
    if (!ch_isrdy(ch)) return 0;    /* la ISRDY */
    *v = ch_read_c(ch);             /* la IN */
    return 1;
}
\end{lstlisting}

Ne discendono due conseguenze. La prima è che l'ospite subisce la stessa
\emph{backpressure} di una cella: se il valore precedente non è ancora stato
consumato, la spinta viene rifiutata e il rifiuto viene contato. La seconda è
che le chiamate vanno effettuate \emph{prima} del passo di clock, nella stessa
finestra in cui le celle calcolano, così che anche l'ospite paghi il ciclo di
latenza per salto.

Alimentazione e drenaggio dei canali di bordo non sono simmetrici. In assenza di
alimentazione, i canali di bordo restano vuoti in modo permanente: è un bordo
aperto da cui non giunge alcun dato, e il programma in attesa gira a vuoto senza
che la griglia si blocchi. In assenza di \textbf{drenaggio}, invece, una
pubblicazione di perimetro resta piena in modo permanente e la cella che intende
pubblicare in quella direzione resta bloccata sulla propria \texttt{SETRDY}. Non
è un difetto, ma la conseguenza diretta della stessa regola che garantisce la
consegna: se nessuno consuma, nessuno può pubblicare. Per questo il simulatore
drena \textbf{sempre}, mentre alimenta solo su richiesta.

\section{Il modello RTL}
\label{sec:rtl}

Il canale e l'interfaccia di cella sono stati riscritti in SystemVerilog. Il
confronto fra le due descrizioni mostra ciò che il passaggio a RTL elimina e
ciò che, al contrario, costa di più.

\subsection{Ciò che il modello RTL elimina}

I campi \texttt{\_next} e la funzione di aggiornamento di fine ciclo
\textbf{non esistono} nel Verilog. Il doppio buffer che in C va emulato a mano
per ottenere la latenza di un ciclo per salto è la semantica nativa
dell'assegnamento non bloccante entro un blocco sincrono:
tutti i registri campionano il valore vecchio e cambiano insieme sul fronte. In
hardware la fedeltà sistolica non ha costo, mentre in C richiede tre campi in
più. Sparisce anche il vincolo sulla \texttt{OUT}, che serviva a impedire due
caricamenti nello stesso ciclo simulato: un ingresso campionato sul fronte non
lo consente per costruzione.

\subsection{Ciò che in RTL costa più che in C}

I registri dati sono due e non uno. Poiché la \texttt{OUT} non è mai rifiutata a
canale pieno, il valore appena caricato e quello pubblicato e non ancora letto
possono differire nello stesso istante, e non possono condividere gli stessi
flip-flop.

I costi, misurati tramite sintesi e non stimati, sono i seguenti:

\begin{center}
\begin{tabular}{lrrr}
    \toprule
     & per canale & per cella & griglia $12 \times 12$ \\
    \midrule
    flip-flop & 66 & 264 & 38\,016 \\
    porte combinatorie & 6 & 24 & 3\,456 \\
    fili verso il vicino & 33 & 132 & --- \\
    fili dal vicino & 1 & 4 & --- \\
    \bottomrule
\end{tabular}
\end{center}

I 66 flip-flop corrispondono ai due registri dati da 32 bit più i due contatori;
le sei porte sono lo XOR e lo XNOR che costituiscono i due comparatori, più
quattro abilitazioni. L'intero protocollo di sincronizzazione su cui è
costruito il progetto costa dunque \textbf{66 flip-flop e 6 porte per canale}, e
fra due celle adiacenti corrono 68 fili, 34 per verso.

Il modello rende visibile anche il motivo per cui \texttt{SETRDY} può essere
rifiutata e \texttt{ISRDY} no. \texttt{ISRDY} non pilota alcuna abilitazione ed
è la lettura diretta dell'uscita di un comparatore, per cui lo zero è una
risposta e non un rifiuto. \texttt{SETRDY} invece pilota l'abilitazione con lo
stesso filo che restituisce, configurando un \emph{test-and-set}.
Da qui si comprende anche perché il test del consumatore debba occupare una
istruzione separata mentre quello del produttore no: \texttt{IN} impegna già
tutti i 32 bit del registro di destinazione con il dato letto, e una istruzione
RISC-V scrive un solo registro.
\clearpage
\section{Le due regioni parallele}
\label{sec:omp-impl}

La parallelizzazione del simulatore si riduce a due direttive, una per fase:

\begin{lstlisting}[caption={Le due regioni parallele del ciclo di clock},label={lst:omp},language=C,float=htbp]
#pragma omp parallel for schedule(static)
for (int i = 0; i < n; i++) {                /* fase compute */
    if (grid -> risc[i].running) {
        execute_step(&grid -> risc[i]);
    }
}
/* barriera implicita a fine ciclo */

#pragma omp parallel for schedule(static)
for (int i = 0; i < n; i++) {                /* fase commit */
    for (int d = 0; d < 4; d++) {
        ch_commit(&grid -> risc[i].out_ch[d]);
    }
}
\end{lstlisting}

La distribuzione è statica perché ogni iterazione ha lo stesso costo (una
istruzione simulata), e decidere a tempo di esecuzione aggiungerebbe soltanto
lavoro. Non compare alcuna sezione critica né alcuna riduzione, perché
\textbf{non sono necessarie}, per il motivo strutturale già esposto nella
sezione~\ref{sec:parallelizzazione}.

La barriera implicita a fine ciclo parallelo coincide con la separazione fra le
due fasi: nessun thread aggiorna finché tutti non hanno terminato il calcolo.
Rimuoverla con la clausola \texttt{nowait} comprometterebbe il determinismo del
risultato. L'aggiornamento dei canali di bordo resta seriale, ed è lecito
perché sono $2(R+C)$ contro $R \cdot C$.